\documentclass[11pt]{article}
\usepackage[utf-8]{inputenc}
\usepackage{amsmath}
\usepackage{amssymb}
\usepackage{graphicx}
\usepackage{xcolor}
\usepackage{fancybox}
\usepackage{hyperref}
\usepackage{longtable}
\usepackage{booktabs}

\title{\textbf{Appendix UNMAPPED_BNK: DOMAIN-BANK-STRATEGIC-DECISIONS} \\ Strategic Bank Decision-Making System (SBDMS) 1.0 \\ Evidence-Based Framework for Global Bank Leadership}
\author{Evidence-Based Framework for Economic and Social Behavior (EBF)}
\date{January 15, 2026}

\begin{document}

\maketitle


% =============================================================================
\begin{tcolorbox}[colback=blue!5!white, colframe=blue!75!black,
    title=Quick Reference: Appendix UNMAPPED_BNK]
\textbf{Appendix:} BNK (UNKNOWN)

\textbf{Core Question:} What are the key research findings in this literature domain?

\textbf{Key Concepts:}
\begin{itemize}[nosep]
\item Literature synthesis and integration
\item Framework application and validation
\item Cross-domain connections
\end{itemize}

\textbf{Cross-References:} See Chapter linkage below
\end{tcolorbox}


\begin{abstract}
This appendix provides a systematic review and integration of key research findings
in the UNKNOWN domain. It synthesizes literature relevant to the Evidence-Based
Framework (EBF) and identifies connections to the 10C CORE architecture.
\end{abstract}


\begin{tcolorbox}[
    colback=orange!5!white,
    colframe=orange!75!black,
    title={\textbf{Appendix Scope: Ziel / In-Scope / Out-of-Scope / Constraints}},
    fonttitle=\bfseries
]

\textbf{Ziel (Objective):}
\begin{itemize}[nosep]
    \item Synthesize key research findings for UNKNOWN integration
\end{itemize}

\textbf{In-Scope:}
\begin{itemize}[nosep]
    \item Literature review and synthesis
    \item Parameter estimates from empirical studies
    \item Framework integration points
\end{itemize}

\textbf{Out-of-Scope:}
\begin{itemize}[nosep]
    \item Original empirical analysis $\rightarrow$ METHOD appendices
    \item Detailed mathematical derivations $\rightarrow$ FORMAL appendices
\end{itemize}

\textbf{Constraints:}
\begin{itemize}[nosep]
    \item Literature selection based on citation impact and relevance
    \item Parameter estimates subject to study conditions
\end{itemize}

\end{tcolorbox}

\begin{center}
\colorbox{lightcyan!30}{\begin{minipage}{0.95\textwidth}
\textbf{Category:} DOMAIN-BANK (Financial Institutional Strategic Decisions) \\
\textbf{Code:} BE \\
\textbf{Version:} 1.0 (SBDMS - Strategic Bank Decision-Making System) \\
\textbf{Status:} EXPERIMENTAL (Framework-level; operationalization pending) \\
\textbf{Confidence:} MEDIUM-HIGH (Theoretical validity high; empirical validation in progress)
\end{minipage}}
\end{center}

% ============================================================================
%%%%%%%%%%%%%%%%%%%%%%%%%%%%%%%%%%%%%%%%%%%%%%%%%%%%%%%%%%%%%%%%%%%%%%%%%%%%%%%
\section{The Fundamental Question}
\label{sec:bnk-fundamental}
%%%%%%%%%%%%%%%%%%%%%%%%%%%%%%%%%%%%%%%%%%%%%%%%%%%%%%%%%%%%%%%%%%%%%%%%%%%%%%%

\begin{tcolorbox}[
    colback=red!5!white,
    colframe=red!75!black,
    title={\textbf{Fundamental Question}}
]
\textbf{How does unknown integrate with and contribute to the
Evidence-Based Framework for understanding behavioral outcomes?}

This appendix addresses the core challenge of formalizing behavioral principles
within a rigorous theoretical framework while maintaining practical applicability.
\end{tcolorbox}


\section{Executive Summary}

This appendix presents the **Strategic Bank Decision-Making System v1.0 (SBDMS)**, a hierarchical framework for modeling strategic decisions across three operational tiers of global financial institutions.

\subsection{Central Insight}

Strategic banking decisions are NOT independent events but are interconnected through:
\begin{enumerate}
\item **Hierarchical Constraints** (Board risk appetite constrains divisional decisions)
\item **Complementarities** (Financial + Strategic dimensions reinforce, Financial + Risk trade off)
\item **Context Dynamics** (Regulatory environment, market volatility, capital constraints)
\item **Learning Feedback** (Outcomes update parameters for next decisions)
\end{enumerate}

\subsection{The Three Critical Decision Types}

\begin{table}[h]
\centering
\begin{tabular}{|l|r|l|l|}
\hline
\textbf{Decision Type} & \textbf{Frequency} & \textbf{Reversibility} & \textbf{Impact} \\
\hline
Market Entry & 2-5x/year & Low (Years) & €€€ \\
Credit Classification & ~Daily (100+) & Medium (Months) & €€€ \\
Capital Allocation & 1-3x/year & Low (Years) & €€€€ \\
\hline
\end{tabular}
\end{table}

\subsection{The Hierarchical Structure}

\begin{equation}
\text{Board} \xrightarrow{\text{Risk Appetite}} \text{Division} \xrightarrow{\text{Tactic}} \text{Operations} \xrightarrow{\text{Deal}} \text{Outcome}
\end{equation}

Each level has its own **decision model with different parameter weights**, but all constrained by the level above.

% ============================================================================
\section{1. Theoretical Foundation: Why Existing Bank Models Fail}

\subsection{1.1 Traditional Approach (Deficient)}

Most banking decision models use **independent silos**:
\begin{itemize}
\item Credit Risk Model (predicts probability of default independently)
\item Market Risk Model (VaR calculation independent of credit)
\item Strategic Planning Model (growth projections independent of actual risk)
\item Compliance Module (KYC checklist, rule-based)
\end{itemize}

\textbf{Problem:} A decision that is ``financially attractive'' AND ``strategically aligned'' BUT ``operationally complex'' gets treated as three separate yes/no decisions, not as an integrated trade-off.

\subsection{1.2 Why EBF Framework is Better}

The EBF approach captures **complementarities and trade-offs** explicitly:

\begin{equation}
U(\text{Decision}) = f(\underbrace{F, S, R, P}_{\text{core dimensions}}, \underbrace{\gamma_{ij}}_{\text{interactions}}, \underbrace{\Psi}_{\text{context}})
\end{equation}

**Key insight:** A decision's utility is NOT the sum of its components but depends on how components interact.

Example:
\begin{itemize}
\item Financial alone (high ROI): Maybe 60\% likely
\item Strategic fit alone (portfolio diversification): Maybe 50\% likely
\item **Financial + Strategic together:** 85\% likely (complementary: $\gamma_{FS} > 0$)
\item **Financial + Operational Complexity together:** 30\% likely (conflict: $\gamma_{FP} < 0$)
\end{itemize}

% ============================================================================
\section{2. The Six Core Decision Dimensions}

### **Dimension 1: Financial (F) – 30-40\% weight depending on level**

\textbf{Definition:} Quantifiable profit impact of the decision.

**Metrics:**
\begin{itemize}
\item ROI/ROAC (Return on Assets/Capital)
\item Net Present Value (NPV) of decision stream
\item Incremental Revenue vs. Incremental Cost
\item Margin Impact (basis points or absolute)
\end{itemize}

**Measurement scale (F ∈ [0,1]):**
\begin{table}[h]
\centering
\begin{tabular}{|c|l|l|}
\hline
$F$ & \textbf{Assessment} & \textbf{Example} \\
\hline
0.90-1.00 & Exceptional financial return & 25%+ ROI, Clear positive NPV \\
0.75-0.89 & Strong financial case & 15-25\% ROI, Solid NPV \\
0.60-0.74 & Acceptable financial return & 8-15\% ROI, Positive but modest \\
0.40-0.59 & Marginal financial return & 3-8\% ROI, NPV near zero \\
0.00-0.39 & Weak/negative financial return & <3\% ROI, Negative NPV \\
\hline
\end{tabular}
\end{table}

### **Dimension 2: Strategic (S) – 20-40\% weight depending on level**

\textbf{Definition:} Alignment with overall bank strategy and competitive positioning.

**Metrics:**
\begin{itemize}
\item Fit with stated geographic/product strategy
\item Competitive positioning impact (market share, capabilities)
\item Synergies with existing businesses
\item Path to future growth
\end{itemize}

**Measurement scale (S ∈ [0,1]):**
\begin{table}[h]
\centering
\begin{tabular}{|c|l|l|}
\hline
$S$ & \textbf{Assessment} & \textbf{Example} \\
\hline
0.90-1.00 & Perfect strategic fit & Core strategy, Direct synergies \\
0.75-0.89 & Strong strategic fit & Aligned with stated priorities \\
0.60-0.74 & Good strategic fit & Supports secondary priorities \\
0.40-0.59 & Weak strategic fit & Tangential to strategy \\
0.00-0.39 & Poor/negative strategic fit & Contradicts stated strategy \\
\hline
\end{tabular}
\end{table}

### **Dimension 3: Risk (R) – 25-35\% weight depending on level**

\textbf{Definition:} Net risk exposure (credit, market, operational, reputational, regulatory).

**Metrics:**
\begin{itemize}
\item Probability of adverse outcome
\item Magnitude of potential loss (worst-case)
\item Correlation with existing portfolio
\item Concentration risk contribution
\end{itemize}

**Measurement scale (R ∈ [0,1], note: INVERSE – high R = bad):**
\begin{table}[h]
\centering
\begin{tabular}{|c|l|l|}
\hline
$R$ & \textbf{Risk Level} & \textbf{Example} \\
\hline
0.90-1.00 & Unacceptable risk & >30\% default probability, regulatory veto \\
0.75-0.89 & High risk & 15-30\% default probability \\
0.60-0.74 & Moderate risk & 5-15\% default probability \\
0.40-0.59 & Manageable risk & 2-5\% default probability \\
0.00-0.39 & Low risk & <2\% default probability, pristine credit \\
\hline
\end{tabular}
\end{table}

**Interpretation:** For risk, we INVERT so that low R = good decision. Alternatively, model as $\text{Risk Acceptability} = (1 - R)$.

### **Dimension 4: Practical (P) – 10-15\% weight depending on level**

\textbf{Definition:} Operational feasibility and execution complexity.

**Metrics:**
\begin{itemize}
\item Implementation timeline (can we execute in required timeframe?)
\item Organizational complexity (how many teams involved?)
\item Data availability (do we have information to make decision?)
\item Technical capability (do we have expertise in-house?)
\end{itemize}

**Measurement scale (P ∈ [0,1]):**
\begin{table}[h]
\centering
\begin{tabular}{|c|l|l|}
\hline
$P$ & \textbf{Feasibility} & \textbf{Example} \\
\hline
0.90-1.00 & Highly feasible & Existing team, rapid execution possible \\
0.75-0.89 & Feasible & Minor resource additions needed \\
0.60-0.74 & Moderately feasible & Some complexity, longer timeline \\
0.40-0.59 & Challenging & Significant operational effort required \\
0.00-0.39 & Unfeasible & Beyond current capability or timeline \\
\hline
\end{tabular}
\end{table}

### **Dimension 5: Emotional (E) – 5-15\% weight**

\textbf{Definition:} Management conviction, team alignment, leadership buy-in.

**Metrics:**
\begin{itemize}
\item CEO/Board commitment level
\item Team morale/enthusiasm (will people execute?)
\item Organizational consensus (coalition building)
\item Risk of post-decision regret or abandonment
\end{itemize}

**Measurement scale (E ∈ [0,1]):**
\begin{table}[h]
\centering
\begin{tabular}{|c|l|l|}
\hline
$E$ & \textbf{Conviction} & \textbf{Example} \\
\hline
0.90-1.00 & Unanimous conviction & CEO + Board + Team fully aligned \\
0.75-0.89 & Strong conviction & Leadership enthusiastic, minor dissent \\
0.60-0.74 & Moderate conviction & Mixed views, manageable consensus \\
0.40-0.59 & Weak conviction & Significant internal debate \\
0.00-0.39 & No conviction & Deep skepticism, likely to be reversed \\
\hline
\end{tabular}
\end{table}

### **Dimension 6: Regulatory (Reg) – 5-10\% weight**

\textbf{Definition:} Regulatory approval probability and compliance sustainability.

**Metrics:**
\begin{itemize}
\item Regulatory approval likelihood (will regulators approve?)
\item Compliance sustainability (can we stay compliant 5+ years?)
\item Reputation risk with regulators
\item Changes to regulatory environment forecast
\end{itemize}

**Measurement scale (Reg ∈ [0,1]):**
\begin{table}[h]
\centering
\begin{tabular}{|c|l|l|}
\hline
$\text{Reg}$ & \textbf{Regulatory Status} & \textbf{Example} \\
\hline
0.90-1.00 & Clear regulatory approval & Pre-approved, aligned with regulations \\
0.75-0.89 & Likely regulatory approval & 80\%+ probability of approval \\
0.60-0.74 & Moderate regulatory fit & 60-80\% approval probability \\
0.40-0.59 & Uncertain regulatory status & 40-60\% approval probability \\
0.00-0.39 & Regulatory barriers & <40\% approval probability, likely to be blocked \\
\hline
\end{tabular}
\end{table}

% ============================================================================
\section{3. Complementarity Parameters: The Interactions ($\gamma$)}

Strategic decisions exhibit **synergies and conflicts** between dimensions.

\subsection{3.1 Synergistic Complementarities ($\gamma > 0$)}

\begin{table}[h]
\centering
\begin{tabular}{|l|r|l|}
\hline
\textbf{Pair} & $\gamma$ & \textbf{Interpretation} \\
\hline
Financial + Strategic & +0.15 & Profit + Alignment = Ideal \\
Risk $\times$ Practical & +0.10 & Low risk + feasible = Low execution risk \\
Strategic + Emotional & +0.12 & Strategy + conviction = Execution success \\
Practical + Regulatory & +0.08 & Feasible + compliant = Sustainable \\
\hline
\end{tabular}
\end{table}

**Mechanism:** When a decision is both profitable AND strategically aligned, synergy emerges because:
- Strategic alignment enables pricing premium (Financial $\uparrow$)
- Profitable decisions attract executive attention and resources
- Result: Combined effect > sum of parts

### **3.2 Conflictual Complementarities ($\gamma < 0$)**

\begin{table}[h]
\centering
\begin{tabular}{|l|r|l|}
\hline
\textbf{Pair} & $\gamma$ & \textbf{Interpretation} \\
\hline
Financial + Risk & -0.20 & High-profit often means high-risk \\
Practical + Risk & -0.15 & Complex execution compounds risk \\
Strategic + Risk & -0.10 & Growth-oriented strategy may increase risk \\
Emotional + Regulatory & -0.12 & Conviction-driven decisions may violate rules \\
\hline
\end{tabular}
\end{table}

**Mechanism:** When a decision offers exceptional returns, the risk is often elevated because:
- Market competition has already priced out safe profits
- Exceptional returns come from untested markets or products
- Result: Strong negative correlation between F and R

% ============================================================================
\section{4. Context Dependencies ($\Psi$)}

Strategic decisions are **not made in vacuum** – context shifts decision-making completely.

\subsection{4.1 Key Context Dimensions}

\begin{table}[h]
\centering
\begin{tabular}{|l|l|l|}
\hline
\textbf{Context Factor} & \textbf{Scale} & \textbf{Impact} \\
\hline
Market Volatility (VIX) & 10-80 & Higher VIX $\rightarrow$ higher risk weight \\
Regulatory Tightness & Loose-Tight & Tight regulation $\rightarrow$ higher Reg weight \\
Capital Position & Surplus-Deficit & Deficit $\rightarrow$ higher Financial weight \\
Risk Appetite & Conservative-Aggressive & Appetite sets constraint on acceptable Risk level \\
Economic Cycle & Expansion-Contraction & Expansion $\rightarrow$ higher Strategic weight \\
\hline
\end{tabular}
\end{table}

\subsection{4.2 Context-Dependent Weighting}

\textbf{Normal Environment:}
$$\text{Weights} = \{F: 0.35, S: 0.25, R: 0.25, P: 0.10, E: 0.05\}$$

\textbf{Crisis Environment (High Volatility):}
$$\text{Weights} = \{F: 0.20, S: 0.15, R: 0.50, P: 0.10, E: 0.05\}$$
(Risk dominates decision-making)

\textbf{Growth Mode (Expansion):}
$$\text{Weights} = \{F: 0.40, S: 0.35, R: 0.15, P: 0.05, E: 0.05\}$$
(Financial and Strategic dominate; risk accepted)

% ============================================================================
\section{5. The Generic Utility Model}

\subsection{5.1 Mathematical Formulation}

The utility of a strategic decision is:

\begin{equation}
U(\text{Decision}) = \beta_0 + \sum_{i} \beta_i C_i + \sum_{i<j} \gamma_{ij} C_i C_j + \beta_\Psi \Psi + \epsilon
\end{equation}

Where:
\begin{align*}
C_i &\in \{F, S, R, P, E, \text{Reg}\} \quad \text{(Core dimensions, each } \in [0,1]) \\
\gamma_{ij} &\in [-0.25, +0.25] \quad \text{(Complementarity parameters)} \\
\Psi &= \{\text{Market Vol, Regulatory, Capital, Risk Appetite}\} \quad \text{(Context)} \\
\beta &= \text{Estimated coefficients (different per hierarchy level)} \\
\epsilon &\sim N(0, \sigma^2) \quad \text{(Unobserved randomness)}
\end{align*}

\subsection{5.2 Decision Probability}

The probability that a decision is **approved** (or should be approved):

\begin{equation}
P(\text{Approve}) = \frac{1}{1 + \exp(-U(\text{Decision}))}
\end{equation}

This **logistic transformation** ensures $P \in [0, 1]$.

% ============================================================================

%%%%%%%%%%%%%%%%%%%%%%%%%%%%%%%%%%%%%%%%%%%%%%%%%%%%%%%%%%%%%%%%%%%%%%%%%%%%%%%
\subsection{Core Axioms}
\label{sec:bnk-axioms}
%%%%%%%%%%%%%%%%%%%%%%%%%%%%%%%%%%%%%%%%%%%%%%%%%%%%%%%%%%%%%%%%%%%%%%%%%%%%%%%

\begin{tcolorbox}[colback=yellow!5!white,colframe=yellow!75!black,title=Axiom UNMAPPED_BNK-1: Fundamental Principle]
\textbf{Axiom UNMAPPED_BNK-1 (Core Assumption):}
The fundamental principle underlying this appendix is that behavioral outcomes
emerge from the interaction of individual characteristics and contextual factors.
\end{tcolorbox}

\begin{tcolorbox}[colback=yellow!5!white,colframe=yellow!75!black,title=Axiom UNMAPPED_BNK-2: Complementarity]
\textbf{Axiom UNMAPPED_BNK-2 (Complementarity):}
Components of the framework exhibit complementarity ($\gamma > 0$), meaning
that combined effects exceed the sum of individual effects.
\end{tcolorbox}



%%%%%%%%%%%%%%%%%%%%%%%%%%%%%%%%%%%%%%%%%%%%%%%%%%%%%%%%%%%%%%%%%%%%%%%%%%%%%%%
\section{Key Results}
\label{sec:bnk-results}
%%%%%%%%%%%%%%%%%%%%%%%%%%%%%%%%%%%%%%%%%%%%%%%%%%%%%%%%%%%%%%%%%%%%%%%%%%%%%%%

The analysis yields the following key results:

\begin{enumerate}
    \item \textbf{Result 1:} [Primary finding with quantitative support]
    \item \textbf{Result 2:} [Secondary finding with implications]
    \item \textbf{Result 3:} [Practical application or boundary condition]
\end{enumerate}


\section{6. Hierarchical Integration}

Strategic banking decisions are constrained by hierarchy.

\subsection{6.1 The Three Decision Levels}

\textbf{LEVEL 1: BOARD OF DIRECTORS}
\begin{itemize}
\item **Decisions:** Risk appetite, M\&A, new business entry, regulatory response
\item **Horizon:** 3-5 years
\item **Parameter emphasis:** Strategic (40\%), Financial (30\%), Risk (30\%)
\item **Constraint mechanism:** Board sets ``Risk Appetite Mandate'' = max acceptable Risk level
\end{itemize}

\textbf{LEVEL 2: BUSINESS DIVISION LEADERSHIP}
\begin{itemize}
\item **Decisions:** Product launches, customer segment focus, capital deployment
\item **Horizon:** 1-2 years
\item **Parameter emphasis:** Financial (35\%), Strategic (25\%), Risk (25\%), Practical (15\%)
\item **Constraint from L1:** ``No decision can exceed Risk Appetite Mandate''
\item **Enablement to L3:** Division provides ``default risk capacity'' to front-line
\end{itemize}

\textbf{LEVEL 3: OPERATIONAL / FRONT-LINE}
\begin{itemize}
\item **Decisions:** Individual customer on-boarding, transaction approval, pricing
\item **Horizon:** Days to weeks
\item **Parameter emphasis:** Risk (35\%), Financial (30\%), Regulatory (20\%), Practical (15\%)
\item **Constraint from L2:** ``No customer exposure can exceed allocated risk budget''
\item **Feedback to L2:** Actual defaults/losses update risk models
\end{itemize}

\subsection{6.2 Constraint Propagation}

\begin{equation}
\text{Board Risk Appetite} \rightarrow \text{Division Risk Budget} \rightarrow \text{RM Risk Allowance per Customer}
\end{equation}

Example: If Board says ``Max 5\% operational risk,'' and Division has 100 RMs with 50 customers each:
- Division allocates: 5\% / 5,000 customers = 0.001\% per-customer maximum
- RM can only approve customers below this threshold

% ============================================================================
\section{7. The Three Strategic Decision Models}

\subsection{7.1 Decision Type 1: Market Entry Decisions}

\textbf{Question:} Should we enter Market X with Product Y for Segment Z?

\textbf{Key Dimensions:}
\begin{itemize}
\item **F (Financial):** Market size estimate, growth rate, penetration probability
\item **S (Strategic):** Geographic diversification, product ecosystem fit, competitive response
\item **R (Risk):** Regulatory barriers, currency risk, geopolitical risk, protectionist response
\item **P (Practical):** Time to profitability, capability gaps, operational readiness
\item **E (Emotional):** Management conviction (Do leaders believe in this market?)
\item **Reg:** Host country regulatory approval, capital requirements, ongoing compliance
\end{itemize}

\textbf{Typical Complementarities:}
\begin{itemize}
\item $\gamma(F, S) = +0.15$: High-growth market that's strategically aligned = exceptional
\item $\gamma(S, R) = -0.15$: Strategically important BUT high-risk = tension
\item $\gamma(P, R) = -0.15$: Complex entry + high risk = compounded execution risk
\end{itemize}

\textbf{Prediction:}
$$P(\text{Entry | Approval}) = \sigma(U_{\text{Market Entry}})$$

If $P > 0.70$: GO (high confidence in approval)
If $0.50 < P < 0.70$: CONDITIONAL (needs additional analysis)
If $P < 0.50$: NO-GO (insufficient support)

\subsection{7.2 Decision Type 2: Credit / Risk Classification}

\textbf{Question:} Should we approve this customer, at what risk rating, and with what conditions?

\textbf{Key Dimensions:}
\begin{itemize}
\item **F (Financial):** Expected spread revenue, probability-of-default, loss-given-default
\item **S (Strategic):** Strategic customer relationship, relationship economics over time
\item **R (Risk):** Credit quality, sector concentration, correlation with portfolio
\item **P (Practical):** Information completeness, monitoring feasibility, reporting capability
\item **Reg (Regulatory):** KYC/AML status, sanctions check, compliance sustainability
\item **E (Emotional):** Relationship manager conviction (Will they support ongoing monitoring?)
\end{itemize}

\textbf{Typical Complementarities:}
\begin{itemize}
\item $\gamma(\text{Reg}, P) = +0.10$: Compliance-ready + full data = operational efficiency
\item $\gamma(F, R) = -0.20$: High-spread customer often has higher default risk
\item $\gamma(F, S) = +0.10$: Profitable relationship + strategic importance = sustainable
\end{itemize}

\textbf{Prediction:}
$$P(\text{Approve | Risk Class K}) = \sigma(U_{\text{Credit}})$$

If $P > 0.80$: APPROVE at standard terms
If $0.60 < P < 0.80$: APPROVE with enhanced conditions (higher spread, lower LTV)
If $0.40 < P < 0.60$: REJECT or refer to credit committee
If $P < 0.40$: HARD REJECT

\subsection{7.3 Decision Type 3: Capital Allocation / M\&A}

\textbf{Question:} Should we acquire this bank/asset or invest in this technology?

\textbf{Key Dimensions:}
\begin{itemize}
\item **F (Financial):** Acquisition price, synergy estimates, NPV calculation, payback period
\item **S (Strategic):** Geographic expansion, capability acquisition, market share gain
\item **R (Risk):** Integration complexity, regulatory approval risk, technology risk
\item **P (Practical):** Time to integration, talent retention, system integration challenges
\item **E (Emotional):** CEO conviction, board alignment, cultural fit perception
\item **Reg:** Regulatory approval probability, capital requirement implications
\end{itemize}

\textbf{Typical Complementarities:}
\begin{itemize}
\item $\gamma(F, S) = +0.20$: Profitable acquisition + strategic fit = exceptional deal
\item $\gamma(R, P) = -0.25$: High integration risk + operational complexity = major problem
\item $\gamma(S, E) = +0.15$: Strategic + leadership conviction = execution success
\end{itemize}

\textbf{Prediction:}
$$P(\text{Proceed to Acquisition}) = \sigma(U_{\text{M\&A}})$$

If $P > 0.75$: GO – proceed to negotiation
If $0.50 < P < 0.75$: CONDITIONAL – needs due diligence on R and P
If $P < 0.50$: NO-GO – abandon or restructure deal

% ============================================================================
\section{8. Parametrization & Calibration}

\subsection{8.1 Parameter Sources (Hybrid Approach)}

**Literature-Based:**
- Default probability models (credit risk literature)
- Strategic fit assessments (business strategy literature)
- Operational complexity scoring (project management)

**Data-Based:**
- Historical outcomes from similar decisions
- Actual customer defaults (refine credit models)
- M\&A success rates (refine acquisition parameters)

**Expert Judgment:**
- Management assessment of risk appetite
- Business unit conviction scores
- Regulatory relationship status

\subsection{8.2 Calibration Example: Credit Decision Model}

For a typical mid-market corporate customer:

\begin{align*}
\beta_0 &= -2.5 \quad \text{(intercept: default expectation)} \\
\beta_F &= 1.2 \quad \text{(financial dimension strongly positive)} \\
\beta_S &= 0.8 \quad \text{(strategic fit moderately positive)} \\
\beta_R &= -2.0 \quad \text{(risk dimension strongly negative)} \\
\beta_P &= 0.6 \quad \text{(practical feasibility moderately positive)} \\
\beta_{\text{Reg}} &= 0.4 \quad \text{(regulatory dimension weakly positive)} \\
\gamma_{F,S} &= +0.15 \quad \text{(synergy between profit and strategy)} \\
\gamma_{F,R} &= -0.20 \quad \text{(trade-off between return and risk)}
\end{align*}

## Example Calculation:

Customer profile:
\begin{align*}
F &= 0.75 \quad \text{(strong financial returns, 12\% spread)} \\
S &= 0.60 \quad \text{(strategic customer, but not core)} \\
R &= 0.40 \quad \text{(acceptable risk, BBB- rated)} \\
P &= 0.85 \quad \text{(full KYC data available)} \\
\text{Reg} &= 0.90 \quad \text{(all compliance checks passed)} \\
\Psi &= \text{Normal environment}
\end{align*}

Utility calculation:

\begin{align*}
U &= -2.5 + 1.2(0.75) + 0.8(0.60) - 2.0(0.40) + 0.6(0.85) + 0.4(0.90) \\
&\quad + 0.15(0.75)(0.60) - 0.20(0.75)(0.40) \\
&= -2.5 + 0.90 + 0.48 - 0.80 + 0.51 + 0.36 + 0.068 - 0.060 \\
&= -0.87
\end{align*}

Approval probability:

\begin{equation}
P(\text{Approve}) = \frac{1}{1 + \exp(0.87)} = \frac{1}{1 + 2.39} = 0.295
\end{equation}

\textbf{Interpretation:} With utility = -0.87, the approval probability is only ~30\%, suggesting either:
- This customer should be REJECTED, OR
- Conditions need improvement (e.g., higher spread to improve F, or better strategic fit)

If RM increases spread to lift $F = 0.85$:
$$U = -0.87 + 0.10(1.2 - 2.0) + 0.10(0.15 - 0.20) = -0.87 - 0.096 - 0.005 = -0.971$$

(Actually gets worse! Because higher spread increases risk correlation.)

Better approach: Improve Strategic fit $S = 0.75$ (core customer):
$$U = -0.87 + 0.15(0.15) = -0.87 + 0.0225 = -0.848$$
$$P(\text{Approve}) = 0.30 \text{ (barely improves)}$$

This illustrates why **single-dimension optimization fails** – must optimize across multiple dimensions with complementarities.

% ============================================================================
\section{9. Learning Loop: From Outcomes to Parameter Updates}

Banking is iterative. Decisions lead to outcomes, which should update our models.

\subsection{9.1 The Feedback Cycle}

\begin{equation}
\text{Decision} \rightarrow \text{Implementation} \rightarrow \text{Outcome} \rightarrow \text{Model Update} \rightarrow \text{Better Decisions}
\end{equation}

**Year 1:** Made 100 credit decisions with model predictions.
**Year 2:** Observed actual defaults, spreads earned, customer dropoff.
**Year 3:** Compare predictions vs outcomes → update $\beta$ and $\gamma$ parameters.
**Year 4:** Use updated model for next 100 decisions.

### **9.2 Example: Credit Model Learning**

\begin{table}[h]
\centering
\begin{tabular}{|l|r|r|r|}
\hline
\textbf{Metric} & \textbf{Predicted} & \textbf{Actual} & \textbf{Action} \\
\hline
Default probability (sample mean) & 3.5\% & 4.2\% & Increase $\beta_R$ weight \\
Approval accuracy & 78\% & 71\% & Add F-R interaction term \\
Strategic customers default rate & 2.1\% & 1.8\% & Strategic impact stronger, increase $\beta_S$ \\
\hline
\end{tabular}
\end{table}

Result: Model becomes **calibrated to your institution**, not generic.

% ============================================================================
\section{10. Known Limitations \& Future Work}

\subsection{10.1 Critical Limitations (v1.0)}

\begin{table}[h]
\centering
\begin{tabular}{|l|l|l|}
\hline
\textbf{Limitation} & \textbf{Severity} & \textbf{Roadmap} \\
\hline
No crisis scenarios modeled & MEDIUM & v1.1: Add binary shock scenarios \\
Political risk underspecified & MEDIUM & v1.1: Add geopolitical variables \\
Behavioral biases not modeled & MEDIUM & v2.0: Add cognitive constraints \\
No real-time calibration & HIGH & v2.0: Implement streaming updates \\
Assumes rational decision-making & MEDIUM & v2.0: Add organizational friction \\
\hline
\end{tabular}
\end{table}

### **10.2 Roadmap**

\textbf{v1.0 (Current):} Framework + 3 decision types, static parameters
\textbf{v1.1 (Q2 2026):} Crisis scenarios, real-time parameter updates
\textbf{v2.0 (Q3-Q4 2026):} Behavioral integrations, organizational constraints
\textbf{v3.0 (2027):} Reinforcement learning, autonomous decision support

% ============================================================================
\section{Conclusion}

SBDMS 1.0 provides a **structured, evidenced-based approach** to strategic banking decisions that:

\begin{enumerate}
\item **Integrates multiple dimensions** (Financial, Strategic, Risk, Practical, Emotional, Regulatory)
\item **Captures complementarities explicitly** (not just additive features)
\item **Respects hierarchy** (Board → Division → Operations)
\item **Learns from outcomes** (continuous parameter refinement)
\item **Provides decision support** (not dogmatic automation)
\end{enumerate}

The framework acknowledges that strategic decisions are fundamentally uncertain. The goal is not to eliminate uncertainty but to **structure it systematically**, so that:
- Decisions are transparent and auditable
- Learning accumulates over time
- Hierarchy is respected
- Trade-offs are explicit

% ============================================================================
\section*{References}

\nocite{bcm_master}

This appendix builds on the Evidence-Based Framework Master Bibliography and incorporates:
\begin{itemize}
\item Credit risk literature (Merton, KMV frameworks)
\item Strategic decision-making (Porter, Rumelt)
\item Behavioral finance (Kahneman, Thaler, Fehr)
\item Bank risk management (Basel III/IV, CCAR frameworks)
\item Organizational decision theory (Simon, March)
\end{itemize}


%%%%%%%%%%%%%%%%%%%%%%%%%%%%%%%%%%%%%%%%%%%%%%%%%%%%%%%%%%%%%%%%%%%%%%%%%%%%%%%

%%%%%%%%%%%%%%%%%%%%%%%%%%%%%%%%%%%%%%%%%%%%%%%%%%%%%%%%%%%%%%%%%%%%%%%%%%%%%%%
\section{Critical Foundations and Objections}
\label{sec:bnk-foundations}
%%%%%%%%%%%%%%%%%%%%%%%%%%%%%%%%%%%%%%%%%%%%%%%%%%%%%%%%%%%%%%%%%%%%%%%%%%%%%%%

\subsection{Objection 1: Scope Limitations}

\textbf{Objection:} The framework presented may have limited applicability
outside the specific domain contexts discussed.

\textbf{Response:} We acknowledge domain-specificity. The framework provides
general principles that should be calibrated to specific contexts. Cross-domain
validation remains an open research question.

\subsection{Objection 2: Measurement Challenges}

\textbf{Objection:} Key constructs may be difficult to measure reliably in
practice.

\textbf{Response:} We recommend using validated scales and instruments where
available, and developing domain-specific proxies where necessary. Sensitivity
analysis should assess robustness to measurement uncertainty.



%%%%%%%%%%%%%%%%%%%%%%%%%%%%%%%%%%%%%%%%%%%%%%%%%%%%%%%%%%%%%%%%%%%%%%%%%%%%%%%
\section{Glossary of Symbols}
\label{sec:bnk-glossary}
%%%%%%%%%%%%%%%%%%%%%%%%%%%%%%%%%%%%%%%%%%%%%%%%%%%%%%%%%%%%%%%%%%%%%%%%%%%%%%%

For comprehensive definitions, see Appendix G (Master Glossary).

\begin{table}[htbp]
\centering
\caption{Key Symbols in Appendix UNMAPPED_BNK}
\begin{tabular}{lll}
\toprule
\textbf{Symbol} & \textbf{Meaning} & \textbf{Reference} \\
\midrule
$\gamma$ & Complementarity parameter & Appendix UNMAPPED_B \\
$\Psi$ & Context vector & Appendix UNMAPPED_V \\
$U$ & Utility function & Chapter 3 \\
\bottomrule
\end{tabular}
\end{table}


\section{References}
\label{sec:references}
%%%%%%%%%%%%%%%%%%%%%%%%%%%%%%%%%%%%%%%%%%%%%%%%%%%%%%%%%%%%%%%%%%%%%%%%%%%%%%%

See master bibliography (\texttt{bcm\_master.bib}) for complete citations.

\nocite{bcm_master}

\end{document}
